\SourceLeaf{284}%
\[
A\int\frac{\cos y}{m\cos\psi}\mathrm{d}x
\]
von \(x=0\) bis \(x=h\) ausgedehnt; es wird aber immer erlaubt sein, darin sowohl \(\cos y\) als \(\cos\psi=1\) zu setzen, und für \(m\) denjenigen Werth, welcher in der Linie \(M\) oder für \(y=0\) gilt, wodurch also das Integral
\[
\begin{aligned}
&=A\int\frac{\mathrm{d}x}{m^0(1+\mu x+\mu'xx+\text{u.\,s.\,w.})}\\
&=\frac{A}{m^0}(h-\tfrac{1}{2}\mu hh+(\tfrac{1}{3}\mu\mu-\tfrac{1}{3}\mu')h^3-\text{u.\,s.\,w.})
\end{aligned}
\]
wird. Es ist immer zureichend, den bis auf die Ordnung \(hh\) damit übereinstimmenden Werth
\[
\frac{Ah}{\BareRoot{m^0m'}}
\]
dafür anzunehmen.
\Article{13}
Die Bestimmung der Grössen \(l^0,l'\) geschieht auf folgende Weise. Es sei \(\chi\) der Winkel, welchen an irgend einer Stelle des Grössten Kreisbogens \(N\) dieser in dem Sinne wachsender \(x\) mit dem Meridian in dem Sinne von Norden nach Süden genommen macht, den Winkel von diesem zu jenem in dem Sinne von der Linken nach der Rechten gezählt; es sei ferner \(S\) die Breite an jener Stelle, \(T\) die Länge von einem beliebigen Meridian an ostwärts gerechnet. Man hat dann daselbst
\[
\begin{aligned}
\mathrm{d}S&=-\cos\chi\mathord{.}\mathrm{d}x+\sin\chi\mathord{.}\mathrm{d}y\\
\mathrm{d}T&=-\frac{\sin\chi}{\cos S}\mathrm{d}x-\frac{\cos\chi}{\cos S}\mathrm{d}y
\end{aligned}
\]
und folglich den partiellen Differentialquotienten
\[
\frac{\mathrm{d}m}{m\mathrm{d}y}=\sin\chi\mathord{.}\frac{\mathrm{d}m}{m\mathrm{d}S}-\frac{\cos\chi}{\cos S}\mathord{.}\frac{\mathrm{d}m}{m\mathrm{d}T}
\]

Da nun bei unserer conformen Uebertragung \(m\) von der Länge unabhängig oder \(\tfrac{\mathrm{d}m}{m\mathrm{d}T}=0\) ist, so wird
\[
l=\sin\chi\mathord{.}\frac{\mathrm{d}m}{m\mathrm{d}S}
\]

Bezeichnet man die Werthe von \(\chi\) in den Punkten \(F\) und \(G\) mit \(V^0\) und \(180^0+V'\) (so dass nach gewöhnlichem Sprachgebrauche \(V^0\) das Azimuth des Grössten Kreisbogens \(FG\) in \(F\), und \(V'\) das Azimuth des Grössten Kreisbogens
